\label{sec:one-chain-many-structures}

\begin{topicbox}{One Chain, Many Structures}
\begin{exposition}
The embeddings assemble into a single chain $\mathbb{N}\hookrightarrow\mathbb{W}\hookrightarrow\mathbb{Z}\hookrightarrow\mathbb{Q}\hookrightarrow\mathbb{R}$. Each arrow is injective, preserves the relevant arithmetic, and reflects order; each step also \emph{adds} structure that the previous system lacked.
\end{exposition}
\begin{exposition}
From $\mathbb{N}$ to $\mathbb{W}$ an additive identity $0$ appears. From $\mathbb{W}$ to $\mathbb{Z}$ additive inverses appear, completing a commutative ring. From $\mathbb{Z}$ to $\mathbb{Q}$ multiplicative inverses of nonzero elements appear, completing an ordered field. From $\mathbb{Q}$ to $\mathbb{R}$ order-completeness appears: the gaps in the rational line are filled. The numerals carry through unchanged by compatibility, but the ambient structure grows at every step. This chain is the bridge from number construction to the analysis of the real line.
\end{exposition}
\end{topicbox}
\begin{remark*}[Exposition]
Read as a ladder of structures: counting ($\mathbb{N}$) $\to$ counting with zero ($\mathbb{W}$) $\to$ signed counting / ring ($\mathbb{Z}$) $\to$ ordered field ($\mathbb{Q}$) $\to$ complete ordered field ($\mathbb{R}$). Only the last rung is order-complete, and that completeness is the property the structure-of-the-real-line chapter turns into metric topology.
\end{remark*}
